\documentclass[11pt]{article}
\usepackage[margin=1in]{geometry}
\usepackage{amsmath,amssymb,amsthm}
\usepackage{booktabs}
\usepackage{microtype}
\usepackage[hidelinks]{hyperref}
\usepackage[T1]{fontenc}
\usepackage{lmodern}
\usepackage{xcolor}
\setlength{\parindent}{0pt}
\setlength{\parskip}{0.55em}
\newtheorem{proposition}{Proposition}
\newtheorem{remark}{Remark}
\newcommand{\Fp}{\mathbf F_p}
\newcommand{\Fseven}{\mathbf F_7}

\title{HOL-01: An exact finite monodromy certificate at $p=7$\\for a Barning--Berggren congruence graph}
\author{Alexandre Couret}
\date{24 September 2026 -- version 1.1.1}

\begin{document}
\maketitle

\begin{abstract}
We give a self-contained exact finite computational certificate for a monodromy phenomenon associated with the classical Barning--Berggren matrices. The principal case is the reduction modulo $49$ to modulo $7$ above $r=(3,4,5)^T$. The root fibre has $49$ points; exact finite closure gives a monodromy group of order $686$, a translation kernel of order $49$, and a dihedral linear image of order $14$. Two explicit positive words, $BABC$ and $BABAA$, generate an order-$343$ subgroup; explicit affine formulas identify this subgroup algebraically with $UT_3(\mathbf F_7)$. The same programs perform finite regression checks at $p=3,5,7,11,13$ and exploratory checks of one higher congruence step at $p=3,5,7$. Those additional computations are not a proof for all odd primes or all levels. Version 1.1.1 deliberately restricts the public claim to what is proved or exactly certified.
\end{abstract}

\section{Setup}
Let
\[
Q(x,y,z)=x^2+y^2-z^2,
\qquad
J=\operatorname{diag}(1,1,-1),
\]
and let
\[
M_1=\begin{pmatrix}-1&2&2\\-2&1&2\\-2&2&3\end{pmatrix},\quad
M_2=\begin{pmatrix}1&2&2\\2&1&2\\2&2&3\end{pmatrix},\quad
M_3=\begin{pmatrix}1&-2&2\\2&-1&2\\2&-2&3\end{pmatrix}.
\]
The certificate first verifies exactly that
\[
M_i^TJM_i=J,
\]
and that the determinants of $(M_1,M_2,M_3)$ are $(1,-1,1)$.  We take the root
\[
r=(3,4,5)^T.
\]

For a modulus $m$, write $V_m$ for the orbit of $r$ modulo $m$.  Computationally we use the symmetric generating set
\[
S=\{M_1^{\pm1},M_2^{\pm1},M_3^{\pm1}\}.
\]
The labelled graph $\Gamma_m$ has vertex set $V_m$ and, for each $s\in S$, a labelled edge $v\to sv$.  Since every $M_i$ is invertible modulo $m$, adding inverse labels does not change the vertex orbit generated by the positive semigroup.

For a prime $p$, reduction modulo $p$ gives a labelled graph morphism
\[
\pi_p:\Gamma_{p^2}\longrightarrow\Gamma_p.
\]
When every fibre has the same size, this is a covering of labelled directed graphs.  The monodromy group used below is the permutation group induced on the fibre over $r$ by loops in $\Gamma_p$.

\section{Exact affine realization of the fibre action}
The bilinear form associated with $Q$ is
\[
\langle u,v\rangle_J=u^T Jv.
\]
Define the tangent plane
\[
T_p=r^\perp=\{u\in\Fp^3:\langle r,u\rangle_J=0\}.
\]
For each tested prime, the certificate directly enumerates $V_{p^2}$ and verifies that the root fibre is exactly
\[
\pi_p^{-1}(r)=\{r+pu:u\in T_p\}.
\]
Intrinsically this is an affine torsor under the two-dimensional vector space $T_p$. The distinguished lift $r$ supplies an origin for the computation; an identification $T_p\simeq\Fp^2$ additionally requires a choice of basis.

If $g$ represents a loop at $r$ modulo $p$, then $g r\equiv r\pmod p$. Writing a fibre point as $r+pu$, one obtains an exact affine action
\[
u\longmapsto A_g u+b_g
\quad\text{on }T_p.
\]
The program constructs a breadth-first spanning tree of $\Gamma_p$, forms the corresponding Schreier loop generators, converts each loop to its affine pair $(A_g,b_g)$, and closes the finite affine group exactly.  No floating-point arithmetic is used.

\section{The certified case $p=7$}
\begin{proposition}[exact finite certificate]
For the reduction $\Gamma_{49}\to\Gamma_7$ above $r=(3,4,5)^T$, the program certifies:
\begin{align*}
|V_7|&=24, & |V_{49}|&=1176,\\
|\pi_7^{-1}(v)|&=49\quad(v\in V_7), & |\operatorname{Mon}_{49/7}|&=686.
\end{align*}
The kernel of the linear-part map on monodromy consists of all $49$ translations of the affine torsor under $T_7$.  The linear image is a dihedral group of order $14$, and the computed affine group contains a zero-translation complement of order $14$.  After choosing a basis of $T_7$, the computed affine monodromy group is therefore represented as
\[
\Fseven^2\rtimes D_{14},
\]
where $D_{14}$ denotes the dihedral group of order $14$.
\end{proposition}

The dihedral identification is not inferred from a group-order lookup: the program finds an element $\rho$ of order $7$ and an involution $\sigma$, verifies
\[
\sigma\rho\sigma^{-1}=\rho^{-1},
\]
and verifies that $\rho$ and $\sigma$ generate all $14$ elements of the linear image.

\section{Two explicit loops and a Heisenberg subgroup}
Set
\[
A=M_3,\qquad B=M_2,\qquad C=M_1.
\]
With column vectors and left action, the positive word $BABC$ has matrix $CBAB$, while $BABAA$ has matrix $AABAB$.  The certificate recomputes
\[
W_1=CBAB=
\begin{pmatrix}
55&120&132\\
24&55&60\\
60&132&145
\end{pmatrix}
\]
and
\[
W_2=AABAB=
\begin{pmatrix}
49&94&106\\
146&287&322\\
154&302&339
\end{pmatrix}.
\]
Their integer actions on the root are
\[
W_1r=(1305,592,1433)^T,
\qquad
W_2r=(1053,3196,3365)^T.
\]
Both return to $r$ modulo $7$, while modulo $49$ their lifts end at
\[
W_1r\equiv(31,4,12)^T,
\qquad
W_2r\equiv(24,11,33)^T.
\]
The two lifted actions do not commute.

Let $x$ and $y$ denote their permutations of the $49$-point fibre.  Each has cycle signature
\[
7^7,
\]
so each has order $7$.  The certificate then closes the subgroup $\langle x,y\rangle$ and obtains
\[
|\langle x,y\rangle|=343.
\]
Writing $z=[x,y]$, it verifies exactly
\[
x^7=y^7=z^7=1,
\qquad
z\ne1,
\qquad
[z,x]=[z,y]=1,
\]
and that the centre of $\langle x,y\rangle$ has order $7$.

\begin{proposition}[explicit Heisenberg identification]
The subgroup generated by the two positive loops $BABC$ and $BABAA$ is isomorphic to $UT_3(\Fseven)$.
\end{proposition}

\begin{proof}
Use the basis $(r,e_1)$ of $T_7$, with $e_1=(1,-2,-1)^T$. An exact calculation gives the affine actions
\[
g_1(s,t)=(s+6t+5,t+3),\qquad g_2(s,t)=(s+2t,t+3).
\]
For $a,b,c\in\Fseven$, put
\[
T(a,b,c)(s,t)=(s+at+c,t+b).
\]
Then
\[
T(a,b,c)T(a',b',c')=T(a+a',b+b',c+c'+ab').
\]
Thus the transformations $T(a,b,c)$ form the standard affine realization of $UT_3(\Fseven)$. Here $g_1=T(6,3,5)$ and $g_2=T(2,3,0)$. The two vectors $(6,3)$ and $(2,3)$ are linearly independent, since their determinant is $12\equiv5\pmod7$, and
\[
[g_1,g_2]=T(0,0,5)
\]
is a non-trivial central translation. Hence the images of $g_1,g_2$ generate the full quotient $\Fseven^2$ modulo the centre and their commutator generates the centre. Therefore they generate all $7^3$ transformations $T(a,b,c)$, so the subgroup is $UT_3(\Fseven)$.
\end{proof}

This Heisenberg subgroup is the index-two preimage of the cyclic order-$7$ rotation subgroup of the dihedral linear image.

\section{Finite regression checks}
The same exact procedure is run at five primes. These rows are finite regression checks only; no universal-prime theorem is inferred from them.

\begin{center}
\begin{tabular}{rrrrrrr}
\toprule
$p$ & $|V_p|$ & $|V_{p^2}|$ & fibre & $|\mathrm{Mon}|$ & linear & rotation preimage\\
\midrule
3  & 4  & 36    & 9   & 54   & 6  & 27\\
5  & 12 & 300   & 25  & 250  & 10 & 125\\
7  & 24 & 1176  & 49  & 686  & 14 & 343\\
11 & 60 & 7260  & 121 & 2662 & 22 & 1331\\
13 & 84 & 14196 & 169 & 4394 & 26 & 2197\\
\bottomrule
\end{tabular}
\end{center}

For each listed prime, the program verifies all of the following: every reduction fibre has $p^2$ points; the root fibre equals $r+pT_p$; the monodromy order is $2p^3$; the translation kernel is all of $\Fp^2$; the linear image has a certified dihedral presentation of order $2p$ and a zero-translation complement; and the preimage of the rotation subgroup has order $p^3$, centre of order $p$, is non-abelian, and has exponent $p$.

\section{Classical context and the open uniform problem}\label{sec:classical}
The matrix and group-theoretic background is classical, but this release does not promote that background into a completed classification theorem for every odd prime or every congruence level.

Let
\[
\varphi(m,n)=(m^2-n^2,2mn,m^2+n^2).
\]
Direct calculation relates the Barning--Berggren matrices to $2\times2$ matrices acting on $(m,n)$; this is part of the standard modular/orthogonal description of the Pythagorean cone. It motivates the tangent-space coordinates used by the certificate.

There is also an elementary affine formula which is valid whenever a word $W$ fixes $r$ modulo an odd prime $p$. Define
\[
\delta_p(W)=\frac{Wr-r}{p}\pmod p.
\]
For a first-order lift $x=r+pu$ one has
\[
W(r+pu)\equiv r+p\bigl(\bar W u+\delta_p(W)\bigr)\pmod{p^2},
\]
and
\[
\delta_p(W_1W_2)=\delta_p(W_1)+\bar W_1\delta_p(W_2).
\]
Thus the problem can be organized as a linear action together with a $1$-cocycle. This formula does not by itself determine the image of the affine representation.

A theorem valid for every odd prime would still require, in the exact orbit/Schreier formulation used here, a uniform proof of the orbit and fibre identification, the relevant congruence-stabilizer image, the kernel of reduction, and the image of the cocycle. A theorem valid for every step $\Gamma_{p^{n+1}}\to\Gamma_{p^n}$ requires a separate level-$n$ argument. These are research targets, not claims of this release.

The program \texttt{finite\_prime\_regression\_check.py} checks the observed finite pattern at $p=3,5,7,11,13$ and performs exploratory exact checks of $\Gamma_{p^3}\to\Gamma_{p^2}$ for $p=3,5,7$. Those computations are deliberately not used as a substitute for the missing uniform proof.

\section{Reproducibility and claim boundary}
The principal certificate is the Python file \texttt{hol01\_certificate.py}.  It depends only on the Python standard library.  A deterministic \texttt{RESULTS.json} file is regenerated and can be checked byte-for-byte by running
\begin{verbatim}
python3 hol01_certificate.py --check RESULTS.json
python3 independent_p7_permutation_check.py
python3 finite_prime_regression_check.py --check FINITE_REGRESSION_RESULTS.json
\end{verbatim}
A second implementation, \texttt{independent\_p7\_permutation\_check.py}, independently realizes the Schreier loops as permutations of the 49 fibre points and recovers monodromy order 686 and the order-343 subgroup without using the affine-coordinate closure routines.  A third program, \texttt{finite\_prime\_regression\_check.py}, performs the stated finite regression and exploratory tower-step checks and writes \texttt{FINITE\_REGRESSION\_RESULTS.json}.  All three programs exit nonzero if an assertion fails.

The claim boundary is strict. Version 1.1.1 certifies the finite $p=7$ statements above, proves the explicit Heisenberg subgroup algebraically, and records the other primes and higher-level checks only as finite evidence. It does not prove the same monodromy group for every odd prime and does not prove stability at every congruence level. No priority is claimed over the classical Barning--Berggren tree, Romik dynamics, graph-covering theory, or standard orthogonal/congruence-group methods.

\section{Literature context}
The tree goes back to Berggren \cite{Berggren1934}, Barning \cite{Barning1963} and Hall \cite{Hall1970}.  Romik's work gives a modern dynamical formulation of the classical Barning tree and explicitly uses the same three matrices \cite{Romik2008}.  Cha and Concei\c{c}\~ao give a recent algebraic generalization of Berggren's theorem and discuss orthogonal-group structure for Pythagorean forms \cite{ChaConceicao2023}.  The affine-sieve and thin-orbit literature places Pythagorean triples in a broader congruence/strong-approximation setting \cite{BGS2010,Ehrman2019}.  A 2026 preprint of Paige further develops one-dimensional dynamics induced from Berggren trees for homogeneous quadratic equations in three variables \cite{Paige2026}.

The passage from $2\times2$ to $3\times3$ matrices used in Section~\ref{sec:classical} is classical: Alperin relates the Pythagorean tree to the modular group \cite{Alperin2005}, and Kontorovich and Oh write the parametrisation of Pythagorean triples as the homomorphism $SL_2(\mathbb{R})\to SO(2,1)$ \cite{KO2012}.  Berggren trees for other ternary quadratic forms are constructed by Cha, Nguyen and Tauber \cite{CNT2018}; the description of graph coverings by permutation voltages goes back to Gross and Tucker \cite{GT1977}; and the reduction of primitive Pythagorean triples modulo odd prime powers, with its counting, is treated in \cite{Biddle2026}.

A targeted literature search carried out on 21 September 2026 did not locate a source stating the exact finite certificate recorded here (in particular the specific $49\to7$ monodromy group together with the two positive loop words generating the order-$343$ subgroup).  This observation is recorded only as search context, not as a legal or historical novelty determination.  The surrounding structures are classical; the present release offers the explicit reproducible finite certificate and does not infer a general novelty claim from the search.

\begin{thebibliography}{99}
\bibitem{Romik2008}
D. Romik,
\emph{The dynamics of Pythagorean triples},
Trans. Amer. Math. Soc. 360 (2008), 6045--6064.
DOI: \href{https://doi.org/10.1090/S0002-9947-08-04467-X}{10.1090/S0002-9947-08-04467-X}.

\bibitem{ChaConceicao2023}
B. Cha and R. Concei\c{c}\~ao,
\emph{A polynomial analogue of Berggren's theorem on Pythagorean triples},
arXiv:2310.02144 (2023).

\bibitem{BGS2010}
J. Bourgain, A. Gamburd and P. Sarnak,
\emph{Affine linear sieve, expanders, and sum-product},
Invent. Math. 179 (2010), 559--644.
DOI: \href{https://doi.org/10.1007/s00222-009-0225-3}{10.1007/s00222-009-0225-3}.

\bibitem{Ehrman2019}
M. Ehrman,
\emph{Almost Primes in Thin Orbits of Pythagorean Triangles},
Int. Math. Res. Not. 2019(11), 3498--3526.
DOI: \href{https://doi.org/10.1093/imrn/rnx191}{10.1093/imrn/rnx191}.

\bibitem{Paige2026}
A. Paige,
\emph{Dynamics of integer zeroes of homogeneous quadratic equations over $\mathbb R^3$},
arXiv:2607.03354 (2026).
\bibitem{Berggren1934}
B. Berggren,
\emph{Pytagoreiska trianglar},
Tidskrift f\"or element\"ar matematik, fysik och kemi 17 (1934), 129--139.

\bibitem{Barning1963}
F. J. M. Barning,
\emph{Over pythagorese en bijna-pythagorese driehoeken en een generatieproces met behulp van unimodulaire matrices},
Math. Centrum Amsterdam, Afd. Zuivere Wisk. ZW-011 (1963).

\bibitem{Hall1970}
A. Hall,
\emph{Genealogy of Pythagorean triads},
Math. Gazette 54 (1970), 377--379.

\bibitem{Alperin2005}
R. C. Alperin,
\emph{The modular tree of Pythagoras},
Amer. Math. Monthly 112 (2005), 807--816.

\bibitem{KO2012}
A. Kontorovich and H. Oh,
\emph{Almost prime Pythagorean triples in thin orbits},
J. Reine Angew. Math. 667 (2012), 89--131.
arXiv:1001.0370.

\bibitem{CNT2018}
B. Cha, E. Nguyen and B. Tauber,
\emph{Quadratic forms and their Berggren trees},
J. Number Theory 185 (2018), 218--256.

\bibitem{GT1977}
J. L. Gross and T. W. Tucker,
\emph{Generating all graph coverings by permutation voltage assignments},
Discrete Math. 18 (1977), 273--283.

\bibitem{Biddle2026}
L. Biddle,
\emph{Primitive Pythagorean triples in Lean and reduction modulo odd prime powers},
undergraduate honors thesis, University of Arkansas (2026).

\end{thebibliography}

\end{document}
